\section{Summary of Contributions to OOD Detection}
\label{sec:summ}

%With increasingly prevalent ML models in the real-world, the problem of designing for open-world settings has become more and more of a priority. OoD detection is a major component in solving this problem but currently struggles with the confounding features present in realistic data. This proposal explores two methods involving the use of foreground and background segmented views to improve discriminative ID class-relevant feature learning and thus OoD detection performance. \sysa uses a view selector to determine which of either the foreground or background view to use for feature fusion and OoD detection. \sysb investigates the benefits of considering all views (including different combinations of fused views) in OoD detection via ensembles of view features. Both methods exhibit significant improvements on challenging datasets over recent state-of-the-art OoD detection baselines, indicating a nascent value in using multiple segmented views for OoD detection. Looking forward, this proposal suggests implementing multi-view foreground-background OoD detection in the CL setting to improve task detection and CIL performance. Furthermore, additional analysis of the features generated by foreground and background views would provide valuable insight into the mechanics of multiple segmented views in OoD detection, leading to better robustness and effective design decisions in applications.
With the increasing prevalence of ML models in real-world settings, the ability of such solutions to correctly handle open-world issues has transitioned from edge case to critical requirement. Modern applications demand that ML systems can easily and reliably operate effectively and efficiently in situations that range from simple (e.g. a bird is not a coyote) to difficult, even for most humans (e.g. coyote versus wolf). With the infinitely large frontier of data and rapid evolution of application scenarios (and thus, demands), preparing any one ML solution for all of the possible issues in realistic open-world domains is intractable. Improving OOD detection, a critical component of open-world ML, helps bridge the gap between ML solutions and real-world applications by enabling ML models to more accurately handle difficult data distributions. One of the major problems in current OOD detection research is exactly the issue of detecting OOD samples that are very similar to ID samples, termed near-OOD detection.

Existing OOD detection methods struggle with near-OOD samples primarily because they do not effectively handle the confounding features present in realistic, near-OOD data. It is paramount to learn and leverage ID class-relevant discriminative features in an efficient and accessible manner. This dissertation proposes two major steps advancing OOD detection towards solving the near-OOD detection task, one leveraging multi-view processing and pre-trained models in the image domain (CASOD), and one proposing an accessible prompt-only method utilizing LLMs in the text domain (ATPO).

In Chapter \ref{diss_casod}, CASOD proposed usage of CVCA feature fusion to intelligently extract discriminative features from ID data using foreground and background views correlated with the original image. Performance on several experimental settings, including the challenging ImageNet-1k dataset against near-OOD datasets of NINCO, SSB-Hard, and Food, demonstrated the efficacy of CASOD in realistic, demanding circumstances. By leveraging pre-trained models and LoRA components, CASOD offered near-OOD detection performance at a fraction of computational resources required by equivalent systems. By creatively leveraging the increasing availability of segmented image datasets and models, CASOD prepares OOD detection for real-world domains.

In Chapter \ref{diss_atpo}, ATPO highlighted the potential of accessible OOD detection methods via prompt-only OOD detection. Through the use of ICL prompts inducing lexical and semantic alignment in LLMs, state-of-the-art near-OOD detection can be achieved both with respect to the only other prompt-only method but also over other inaccessible OOD detection methods that require LLM features and/or logits (hence requiring expert knowledge and sophisticated ML programming skills). Thorough experiments across near-OOD datasets, including the challenging Banking intent dataset and 20NewsGroups classification dataset, verify the abilities of ATPO. The ease-of-use and near-OOD detection performance of ATPO fundamentally promotes it as a solution for modern open-world applications where such accessibility and effectiveness is in demand.

\section{Advancing OOD Detection Further}
\label{sec:futwork}

\subsection{Multi-View OOD Detection in Continual Learning}
\label{sec:future_work-mvoodcl}

%(connect ood to cl again)

%It was previously discussed how OoD detection is an important component of preparing ML models for use in realistic, open-world settings (see Section \ref{sec:intro}). A promising next step towards fully functional open-world ML designs is incorporating OoD detection capabilities in the \emph{continual learning} (CL) paradigm. CL aims to incrementally learn a sequence of tasks \cite{thrun1995lifelong,hadsell2020embracing,chen2018lifelong}. This problem closely mirrors how humans learn in the real-world: tasks are continuously encountered in life and humans constantly adjust to solve each of these tasks without forgetting how to complete any previous task. The ability to understand when a new task has arrived is critical to CL and learning in the open-world setting \cite{kim2022theoretical}.
A promising next-step towards open-world ML design involves incorporating OOD detection in the \textit{continual learning} (CL) paradigm. CL learns a sequence of tasks \cite{de2021continual,kim2022theoretical} and mimics how humans learn in the real-world: tasks are continuously encountered and humans adjust to solve any given new task, all while maintaining any required knowledge for previously-encountered tasks. The ability to understand when a new task has been encountered is critical to CL, and learning in general in the open-world setting, and is fundamentally tied to OOD detection \cite{kim2022theoretical}.

%(define CIL and why it's important, Kim's work on ood directly improving cl)

%It has been shown that OoD detection is a fundamental component in CL, with improved OoD detection directly improving CL Class-Incremental Learning (CIL) performance \cite{kim2022theoretical}. CIL is a form of CL setting that learns a sequence of tasks $1,2,...,T$. Each task $k$ has a training dataset $\mathcal{D}_{k} = \{(x^{i}_{k}, y^{i}_{k})^{n_{k}}_{i=1}\}$ with $n_{k}$ as the number of data samples in task $k$. $x^{i}_{k} \in X$ are the input data samples, and $y^{i}_{i} \in Y_{k}$ are the class labels corresponding to each input data sample. Note that $X$ is the set of all data samples and $Y_{k}$ is the set of all classes in task $k$. In this problem definition, it is assumed that all $Y_{k}$'s are disjoint i.e. $Y_{k} \bigcap Y_{k'} = \emptyset, \forall k \neq k'$ and $\bigcup^{T}_{k=1} Y_{k} = Y$. CIL aims to construct a single predictive function or classifier $f: X \mapsto Y$ that classifies each given test instance $x'$ to its class label $y$ \cite{kim2022theoretical}. In CIL, the model is never provided with any indication of which task an encountered sample belongs to. Therefore, a CIL model must \emph{detect} when test instances do not belong to any of the previously encountered ID classes. This requirement is exactly the OoD detection problem. In fact, good classification performance within a given task and good OoD detection performance are necessary and sufficient conditions for good CIL performance \cite{kim2022theoretical}.

%(explain how mv ood can be included in this framework and the expected results)

%The proposed works (see Section \ref{sec:propwork}) indicate that the use of multiple views, specifically foreground-background segmented views, leads to improved OoD detection performance. By applying procedures such as \sysa or \sysb in a CL setting, it is expected that CIL performance will drastically improve when multi-view OoD detection methods are applied with CL models that guarantee within-task classification performance. An initial experiment to validate this idea would use a state-of-the-art CL method, such as HAT \cite{serra2018overcoming}, on an MVCNN architecture. The HAT-MVCNN model would then learn features and logits for OoD detection and ID classification. Extensive experimentation on CL datasets, such as CIFAR-10 \cite{krizhevsky2009learning} divided into 5 tasks or Tiny ImageNet \cite{le2015tiny} divided into 10 tasks, would be used for evaluation of the proposed idea.
By improving OOD detection, CL performance can also be improved as automatic detection and separation of tasks is a fundamental component of CL \cite{kim2022theoretical}. Therefore, it is natural to consider the success of CASOD and related multi-view approaches applied in the CL domain. Investigation into the integration of foreground-background segmentation and multi-view learning in CL could lead to improvement not only in the OOD detection component of CL but also the classification component of CL, introducing multiple forms of improvement. 

\subsection{Multi-View OOD Detection with Semantic Segmentation}
\label{sec:future_work-semseg}

CASOD, while effective, is limited to handling exactly two views in addition to the original image. More specifically, CASOD is designed to take as input the original view and foreground and background segmented views from the original image. However, the space of views is significantly larger and more variable than only the foreground and background of an image. A natural consideration is the potential improvement from leveraging the \textit{semantic segmentation} of an image as additional views for OOD detection \cite{guo2018review}. That is, given an image, variable amount of semantic objects can be segmented from the image. The information in these segmented objects could be beneficial to OOD detection as they each include some amount of ID-specific, correlated information. This perspective is interesting and challenging for multiple reasons. Most significant is the expansion of ATPO to an unknown number of additional views per image. Another challenge is the adaptive fusion of view features with the original view feature (or with each other), as different views may have different semantic contributions with respect to the overall sample. Accounting for such an expansion on the direction of multi-view OOD detection may greatly improve performance in such fine-grained domains as near-OOD detection and is thus of great interest as a future research direction.

\subsection{LLM integration of OOD Detection for Improved Generation}
\label{sec:future_work-llm}

With the continued expansion of LLMs in modern technology and the advent of agentic AI solutions \cite{yuksel2025multi}, improvement of LLM performance has become critical to ensure not only the effective use of tools and products but also to establish reliability and trust in ML-based interfaces. ATPO offers a novel solution to effective OOD detection in a medium that is highly accessible for non-ML-expert users. Integrating ATPO with technologies that utilize LLMs, e.g. agents, could enhance the abilities of such solutions to detect subtle errors in LLM processing of inputs at a low cost to developers (i.e. requiring only a prompt). Investigating the performance effects of including ATPO in such systems could be revealing regarding the intricacies involved in integrating accessible near-OOD detection methods in modular, LLM-based technologies.
